\par\SourceLeaf{425}%
\noindent\hspace*{1.2em}Wir nehmen an, dass die Punkte so geordnet sind, dass sie das von ihnen eingeschlossene Dreieck im selben Sinn umlaufen wie die Punkte $(1),(2),(3)$. Ferner sei $\lambda$ derjenige Pol des Großkreises $L'L''$, der auf derselben Seite liegt wie $L$. Aus unserem vorigen Lemma ergibt sich dann
\[
\begin{aligned}
y'z''-z'y''&=\sin L'L''.\cos\lambda(1)\\
z'x''-x'z''&=\sin L'L''.\cos\lambda(2)\\
x'y''-y'x''&=\sin L'L''.\cos\lambda(3).
\end{aligned}
\]
Multipliziert man also diese Gleichungen mit $x,y,z$ und addiert die Produkte, so erhält man
\[
xy'z''+x'y''z+x''yz'-xy''z'-x'yz''-x''y'z=\sin L'L''.\cos\lambda L,
\]
wofür man nach den bekannten Regeln der sphärischen Trigonometrie auch
\[
\begin{aligned}
&\sin L'L''.\sin LL''.\sin L'\\
&=\sin L'L''.\sin LL'.\sin L''\\
&=\sin L'L''.\sin L'L''.\sin L
\end{aligned}
\]
schreiben kann, wenn $L,L',L''$ die drei Winkel des sphärischen Dreiecks $LL'L''$ bezeichnen. Man überzeugt sich zugleich leicht, dass diese Größe der sechste Teil der Pyramide ist, die zwischen dem Kugelmittelpunkt und den Punkten $L,L',L''$ gebildet wird (und zwar \textls[100]{positiv}, wenn usw.).

\Needspace{5\baselineskip}
\begin{center}8.\end{center}
\noindent\hspace*{1.2em}Die Natur einer krummen Fläche wird durch eine Gleichung zwischen den Koordinaten ihrer Punkte dargestellt, die wir als
\[
f(x,y,z)=0
\]
schreiben. Das vollständige Differential von $f(x,y,z)$ sei
\[
=Pdx+Qdy+Rdz,
\]
wobei $P,Q,R$ Funktionen von $x,y,z$ sind. Wir werden an der Fläche stets zwei Seiten unterscheiden, von denen wir die eine die obere und die andere die untere nennen wollen. Im Allgemeinen wird beim Durchgang durch\space%
